\begin{proof}[Proof of \cref{obj:prop:identification}]
Fix \(x\in\mathcal X\). Write
\[
p'_x=P'(X=x),
\qquad
\mu'_x(a)=\frac{E_{P'}[\mathbf1\{X=x\}Y(a)]}{p'_x},
\qquad a\in[0,1].
\]
By \cref{obj:ass:stratum-overlap}, both \(p_x\) and \(p'_x\) are strictly positive.

\begin{enumerate}
\item Define the centered observed dose by
\[
V=W-a_0=A-a_0+\sigma Z.
\]
For every Borel set \(B\subseteq\mathbb R\), define
\[
\begin{aligned}
\nu_{0x}(B)&=P(A-a_0\in B\mid X=x),
&
\nu'_{0x}(B)&=P'(A-a_0\in B\mid X=x),\\
M_{0x}(B)&=\frac{P(X=x,V\in B)}{p_x},
&
M'_{0x}(B)&=\frac{P'(X=x,V\in B)}{p'_x}.
\end{aligned}
\]
Equality of the observed laws equates their normalized restrictions to each stratum, and therefore
\[
M_{0x}=M'_{0x}.
\]

Let the Gaussian noise measure be
\[
\Phi_\sigma=\mathcal L(\sigma Z),
\qquad
\Phi_0=\delta_0,
\qquad
\delta_0(B)=\mathbf1\{0\in B\}.
\]
By \cref{obj:ass:gaussian-channel}, this is the centered Gaussian probability measure with variance \(\sigma^2\) under either structural law. For Borel sets \(B_Z,B_t\subseteq\mathbb R\), \cref{obj:ass:error-independence} gives
\[
\begin{aligned}
P(Z\in B_Z,A-a_0\in B_t\mid X=x)
&=\frac{P(Z\in B_Z)\,P(X=x,A-a_0\in B_t)}{p_x}\\
&=P(Z\in B_Z)\,\nu_{0x}(B_t).
\end{aligned}
\]
Scaling the error coordinate consequently yields
\[
\mathcal L_P(A-a_0,\sigma Z\mid X=x)
=\nu_{0x}\otimes\Phi_\sigma.
\]
The same factorization holds under \(P'\). Pushing these product measures forward under addition gives
\[
\begin{aligned}
M_{0x}(B)
&=\int_{\mathbb R}\int_{\mathbb R}
\mathbf1\{t+z\in B\}\,\Phi_\sigma(dz)\,\nu_{0x}(dt)
=(\nu_{0x}*\Phi_\sigma)(B),\\
M'_{0x}(B)&=(\nu'_{0x}*\Phi_\sigma)(B).
\end{aligned}
\]

The density representation in \cref{obj:ass:weak-design}, together with the support condition in \cref{obj:ass:latent-support} and \(a_0=1/2\), gives
\[
\nu_{0x}(B)=\int_{B\cap[-1/2,1/2]}g_x(t)\,dt.
\]
Both latent laws are probability measures, and the same support conclusion holds under \(P'\):
\[
\nu_{0x}(\mathbb R)=\nu'_{0x}(\mathbb R)=1,
\qquad
\nu_{0x}\bigl(\mathbb R\setminus[-1/2,1/2]\bigr)
=\nu'_{0x}\bigl(\mathbb R\setminus[-1/2,1/2]\bigr)=0.
\]
Thus
\[
\nu_{0x}*\Phi_\sigma=\nu'_{0x}*\Phi_\sigma
\quad\Longrightarrow\quad
\nu_{0x}=\nu'_{0x},
\]
by \cref{obj:lem:compact-gaussian-convolution-injective}, taking its two auxiliary measures to be zero. Its finiteness, support, and error-scale hypotheses are satisfied, including at \(\sigma=0\).
% lean: CausalSmith.Stat.NoisydoseWeakdesignTransition.latentDoseLaw_eq_of_obsLaw_eq

\item Since \(\beta>0\), the Hölder means are continuous on \([0,1]\). Define, for every Borel \(B\subseteq\mathbb R\),
\[
\begin{aligned}
\nu_{1x}(B)
&=\int_{B\cap[-1/2,1/2]}
\max\{\mu_x(t+a_0),0\}\,\nu_{0x}(dt),\\
\nu'_{1x}(B)
&=\int_{B\cap[-1/2,1/2]}
\max\{\mu'_x(t+a_0),0\}\,\nu'_{0x}(dt).
\end{aligned}
\]
The support of the unmarked laws makes these the marked latent measures in the statement. On this support, \cref{obj:ass:mean-range} makes the positive parts equal to the means themselves. Continuity and boundedness ensure integrability, and hence
\[
\begin{aligned}
\nu_{1x}(\mathbb R)
&=\int_{[-1/2,1/2]}\mu_x(t+a_0)\,\nu_{0x}(dt)\le1,
&
\nu'_{1x}(\mathbb R)&\le1,\\
\nu_{1x}\bigl(\mathbb R\setminus[-1/2,1/2]\bigr)
&=0,
&
\nu'_{1x}\bigl(\mathbb R\setminus[-1/2,1/2]\bigr)&=0.
\end{aligned}
\]

By \cref{obj:lem:realized-dose-mean},
\[
0\le Y\le1
\qquad P\text{-almost surely and }P'\text{-almost surely}.
\]
We may therefore define the finite nonnegative observed marked measures by
\[
M_{1x}(B)=E_P[\mathbf1\{V\in B\}Y\mid X=x],
\qquad
M'_{1x}(B)=E_{P'}[\mathbf1\{V\in B\}Y\mid X=x].
\]
They are determined by the normalized observed laws, so
\[
M_{1x}=M'_{1x}.
\]

For a fixed Borel \(B\), use the bounded measurable test
\[
(a,x',z)\longmapsto
\mathbf1\{x'=x,\ a+\sigma z-a_0\in B\},
\qquad
(a,x',z)\in\mathbb R\times\mathcal X\times\mathbb R.
\]
The bounded-test identity in \cref{obj:lem:realized-dose-mean} gives
\[
\begin{aligned}
M_{1x}(B)
&=\frac{E_P[\mathbf1\{X=x,V\in B\}Y]}{p_x}\\
&=\frac{E_P[\mathbf1\{X=x,V\in B\}\mu_X(A)]}{p_x}\\
&=E_P[\mathbf1\{V\in B\}\mu_x(A)\mid X=x].
\end{aligned}
\]
Using the conditional product law established above, and the supported latent mean, we obtain
\[
\begin{aligned}
M_{1x}(B)
&=\int_{[-1/2,1/2]}\int_{\mathbb R}
\mathbf1\{t+z\in B\}\mu_x(t+a_0)\,
\Phi_\sigma(dz)\,\nu_{0x}(dt)\\
&=(\nu_{1x}*\Phi_\sigma)(B).
\end{aligned}
\]
The same calculation under \(P'\) yields
\[
M'_{1x}=\nu'_{1x}*\Phi_\sigma.
\]
Consequently,
\[
\nu_{1x}*\Phi_\sigma=\nu'_{1x}*\Phi_\sigma
\quad\Longrightarrow\quad
\nu_{1x}=\nu'_{1x},
\]
again by \cref{obj:lem:compact-gaussian-convolution-injective} with the two auxiliary measures zero. The displayed finiteness and support properties discharge its measure hypotheses.
% lean: CausalSmith.Stat.NoisydoseWeakdesignTransition.latentMeanMeasure_eq_of_obsLaw_eq

\item For every integer \(j\ge0\), put
\[
\varepsilon_j=\frac1{j+3},
\qquad
B_j=(-\varepsilon_j,\varepsilon_j).
\]
Then
\[
0<\varepsilon_j\le\frac12,
\qquad
B_j\subseteq[-1/2,1/2].
\]
The lower envelope in \cref{obj:ass:weak-design} implies
\[
g_x(t)\ge c_{\mathrm{lo}}|t|^\kappa>0
\qquad
\text{for Lebesgue-almost every }t\in B_j\setminus\{0\}.
\]
The singleton \(\{0\}\) has Lebesgue measure zero, whereas
\[
\int_{B_j}1\,dt=2\varepsilon_j>0.
\]
If the integral of the nonnegative density over \(B_j\) were zero, the density would vanish almost everywhere there, contradicting this positivity. Thus
\[
\nu_{0x}(B_j)=\int_{B_j}g_x(t)\,dt>0.
\]
This proves every asserted shrinking-neighborhood mass is strictly positive.
% lean: CausalSmith.Stat.NoisydoseWeakdesignTransition.latentDoseLaw_local_mass_pos

\item For \(t\in B_j\), both \(t+a_0\) and \(a_0\) belong to \([0,1]\). The radius-one Hölder bound in \cref{obj:ass:holder-mean} gives
\[
|\mu_x(t+a_0)-\mu_x(a_0)|
\le |t|^\beta
\le\varepsilon_j^\beta.
\]
Integrating this pointwise bound against the latent probability law yields
\[
\left|
\int_{B_j}\bigl(\mu_x(t+a_0)-\mu_x(a_0)\bigr)\,\nu_{0x}(dt)
\right|
\le\varepsilon_j^\beta\,\nu_{0x}(B_j).
\]
Since the marked mass equals the integral of the mean, and the denominator is positive,
\[
\begin{aligned}
\left|
\frac{\nu_{1x}(B_j)}{\nu_{0x}(B_j)}-\mu_x(a_0)
\right|
&=
\frac{
\left|\displaystyle\int_{B_j}
\bigl(\mu_x(t+a_0)-\mu_x(a_0)\bigr)\,\nu_{0x}(dt)\right|
}{\nu_{0x}(B_j)}\\
&\le\varepsilon_j^\beta.
\end{aligned}
\]
% lean: CausalSmith.Stat.NoisydoseWeakdesignTransition.latentLocalRatio_error
Finally,
\[
\varepsilon_j\longrightarrow0,
\qquad
\varepsilon_j^\beta\longrightarrow0,
\]
because \(\beta>0\). Squeezing the nonnegative absolute error therefore proves
\[
\lim_{j\to\infty}\frac{\nu_{1x}(B_j)}{\nu_{0x}(B_j)}
=\mu_x(a_0).
\]
% lean: CausalSmith.Stat.NoisydoseWeakdesignTransition.latentLocalRatio_tendsto

\item Projecting the common observed law onto its stratum coordinate gives
\[
p_x=P(X=x)=P'(X=x)=p'_x.
\]
% lean: CausalSmith.Stat.NoisydoseWeakdesignTransition.strataProb_eq_of_obsLaw_eq

\item The equalities of the unmarked and marked latent measures imply, by uniqueness of densities with respect to their common unmarked probability measure,
\[
\max\{\mu_x(t+a_0),0\}
=\max\{\mu'_x(t+a_0),0\}
\qquad
\text{for }\nu_{0x}\text{-almost every }t\in[-1/2,1/2].
\]
The mean-range bounds remove the positive parts, giving
\[
\mu_x(t+a_0)=\mu'_x(t+a_0)
\qquad
\text{for }\nu_{0x}\text{-almost every }t\in[-1/2,1/2].
\]

The lower design envelope makes \(g_x(t)>0\) for Lebesgue-almost every \(t\) in the full latent interval: the possible exceptional point \(t=0\) is null. Thus, for every Borel \(B\subseteq\mathbb R\),
\[
\nu_{0x}(B)=\int_{B\cap[-1/2,1/2]}g_x(t)\,dt=0
\quad\Longrightarrow\quad
\int_{B\cap[-1/2,1/2]}1\,dt=0.
\]
Indeed, a zero integral of the nonnegative density forces it to vanish almost everywhere on that set, and its almost-everywhere positivity forces the set to be Lebesgue-null. The mean equality consequently transfers to Lebesgue-almost every centered dose:
\[
\mu_x(t+a_0)=\mu'_x(t+a_0)
\qquad
\text{for Lebesgue-almost every }t\in[-1/2,1/2].
\]

For \(t,t'\in[-1/2,1/2]\), the two Hölder bounds give
\[
\begin{aligned}
|\mu_x(t+a_0)-\mu_x(t'+a_0)|&\le|t-t'|^\beta,\\
|\mu'_x(t+a_0)-\mu'_x(t'+a_0)|&\le|t-t'|^\beta.
\end{aligned}
\]
Since \(\beta>0\), both centered mean functions are continuous on the closed interval. If they differed at any point, continuity would preserve their difference on a relative neighborhood of that point. Every such neighborhood has positive Lebesgue measure, including at either endpoint, contradicting their almost-everywhere agreement. Hence
\[
\mu_x(t+a_0)=\mu'_x(t+a_0)
\qquad
\text{for every }t\in[-1/2,1/2].
\]
For every \(a\in[0,1]\),
\[
a-a_0\in[-1/2,1/2],
\qquad
(a-a_0)+a_0=a.
\]
Substitution therefore proves
\[
\mu_x(a)=\mu'_x(a)
\qquad
\text{for every }a\in[0,1].
\]
Since \(x\) was arbitrary, all the preceding conclusions hold in both strata.
% lean: CausalSmith.Stat.NoisydoseWeakdesignTransition.condPotMean_eqOn_of_latent_measures_eq

\item Bounded potential outcomes ensure integrability at \(a_0\). The two strata partition the structural space, so pointwise
\[
Y(a_0)=\sum_{x\in\mathcal X}\mathbf1\{X=x\}Y(a_0).
\]
Taking expectations and using \cref{obj:def:conditional-potential-mean}, with the positive stratum probabilities, gives
\[
\begin{aligned}
\theta(P)
&=E_P[Y(a_0)]\\
&=\sum_{x\in\mathcal X}E_P[\mathbf1\{X=x\}Y(a_0)]\\
&=\sum_{x\in\mathcal X}p_x\mu_x(a_0).
\end{aligned}
\]
The same decomposition applies under \(P'\). Since \(a_0=1/2\in[0,1]\), the identified stratum probabilities and mean functions yield
\[
\theta(P')
=\sum_{x\in\mathcal X}p'_x\mu'_x(a_0)
=\sum_{x\in\mathcal X}p_x\mu_x(a_0)
=\theta(P).
\]
% lean: CausalSmith.Stat.NoisydoseWeakdesignTransition.causalTarget_eq_stratum_sum
% lean: CausalSmith.Stat.NoisydoseWeakdesignTransition.identification
\end{enumerate}
\end{proof}
